\documentclass[12pt, a0paper]{article}

% ===== 页面 =====
\usepackage[
    paperwidth=84.1cm,
    paperheight=118.9cm,
    margin=1.5cm,
    columnsep=1.5cm
]{geometry}

% ===== 中文 =====
\usepackage[fontset=fandol]{ctex}

% ===== 双栏 =====
\usepackage{multicol}

% ===== 表格 =====
\usepackage[table]{xcolor}
\usepackage{colortbl}
\usepackage{multirow}
\usepackage{booktabs}

% ===== 其他 =====
\usepackage{tikz}
\usepackage{enumitem}
\usepackage{amsmath}
\usepackage{amssymb}

% ===== 颜色 =====
\definecolor{myblue}{RGB}{25,55,109}
\definecolor{myorange}{RGB}{200,100,20}
\definecolor{mygreen}{RGB}{30,120,70}
\definecolor{myred}{RGB}{180,40,40}
\definecolor{mygray}{RGB}{130,130,130}
\definecolor{mypurple}{RGB}{100,40,160}
\definecolor{lightblue}{RGB}{235,245,255}

% ===== 字体缩放 (A0) =====
\renewcommand{\normalsize}{\fontsize{20}{27}\selectfont}
\renewcommand{\large}{\fontsize{24}{32}\selectfont}
\renewcommand{\Large}{\fontsize{28}{36}\selectfont}
\renewcommand{\LARGE}{\fontsize{34}{42}\selectfont}
\renewcommand{\huge}{\fontsize{40}{48}\selectfont}
\renewcommand{\Huge}{\fontsize{48}{56}\selectfont}
\renewcommand{\small}{\fontsize{16}{22}\selectfont}
\renewcommand{\footnotesize}{\fontsize{14}{18}\selectfont}

% ===== 收紧间距 =====
\setlength{\parskip}{0.3em}
\setlength{\parindent}{0pt}
\setlength{\multicolsep}{0.5em}
\raggedbottom

% ===== 节标题 =====
\usepackage{titlesec}
\titleformat{\section}[hang]
    {\LARGE\bfseries\color{myblue}}
    {\thesection}
    {0.8em}
    {}
    [\vspace{-1mm}{\color{myblue}\rule{\columnwidth}{1.2pt}}]
\titlespacing{\section}{0pt}{0.3em}{0.2em}

\titleformat{\subsection}[hang]
    {\large\bfseries\color{myblue!80}}
    {\thesubsection}
    {0.8em}
    {}
\titlespacing{\subsection}{0pt}{0.2em}{0.05em}

% ===== 列表 =====
\setlist{nosep, leftmargin=2.5em, itemsep=0.3em}

% ===== 占位图 =====
\newcommand{\placeholder}[2]{%
    \begin{center}
    \colorbox{mygray!12}{\parbox{0.9\columnwidth}{%
        \centering
        \vspace{1.5em}
        {\Large\bfseries\color{mygray} [ 占位图： #1 ]}
        \vspace{0.8em}

        {\small\color{mygray} #2}
        \vspace{1.5em}
    }}
    \end{center}
}

\newcommand{\placeholderh}[3]{%
    \begin{center}
    \colorbox{mygray!12}{\parbox{0.9\columnwidth}{%
        \centering
        \vspace{1.2em}
        {\Large\bfseries\color{mygray} [ 占位图： #1 ]}
        \vspace{0.8em}
        {\small\color{mygray} #2}
        \vspace{#3}
    }}
    \end{center}
}

\newcommand{\modelsizecompare}{%
    \begin{center}
    \colorbox{mygray!8}{\parbox{0.95\columnwidth}{%
        \vspace{0.4em}
        {\large\bfseries\color{myblue}模型规模对比（参数量/计算量）}\par
        \vspace{0.4em}
        \begin{tikzpicture}[x=1cm,y=1cm]
            \fill[mygreen!30] (0,0) rectangle (1.0,0.5);
            \node[anchor=west] at (1.2,0.25) {\normalsize EEGNet \textcolor{mygray}{(示例：0.02M)}};

            \fill[myorange!30] (0,-0.8) rectangle (2.5,-0.15);
            \node[anchor=west] at (2.7,-0.475) {\normalsize EEG Conformer \textcolor{mygray}{(示例：1.5M)}};

            \fill[mypurple!25] (0,-1.7) rectangle (5.0,-0.85);
            \node[anchor=west] at (5.2,-1.275) {\normalsize CBraMod \textcolor{mygray}{(示例：20M)}};
        \end{tikzpicture}
        \vspace{0.3em}\par
        {\footnotesize\color{mygray}注：替换为真实统计值后可表达实际大小对比。}\par
        \vspace{0.4em}
    }}
    \end{center}
}

\begin{document}

% =====================
% 标题
% =====================
\begin{center}
    {\fontsize{48}{58}\selectfont\bfseries\color{myblue}
        EEG 解码模型的发展：\\[2mm]
        从 EEGNet 到 EEG Foundation Model 的复现与分析\\[4mm]
    }

    {\fontsize{24}{32}\selectfont\color{mygray}
        基于 EEGNet、EEG Conformer 和 CBraMod 的比较研究\\[4mm]
    }

    {\fontsize{22}{28}\selectfont
        第一作者 (SUSTech ID) \hspace{2em} 第二作者 (ID) \hspace{2em} 第三作者 (ID)\\[2mm]
    }

    {\fontsize{16}{22}\selectfont\color{mygray}
        南方科技大学生物医学工程系 \hspace{3em}
        BMEB215 机器学习与医学工程应用 · 2026 Spring
    }

    \vspace{3mm}
    {\color{myblue}\hrule height 2pt}
    \vspace{3mm}
\end{center}

\begin{multicols}{2}

% ================================================================
% 1. 研究背景
% ================================================================
\section{研究背景 (Background)}

\subsection{EEG 解码的重要应用}

脑机接口 (BCI) · 睡眠分期 (Sleep Staging) · 神经疾病辅助诊断 · 认知状态监测

\subsection{EEG 解码面临的挑战}

\begin{itemize}
    \item 数据标注成本高 —— 临床、睡眠等任务需要专家标注
    \item 被试间差异显著 —— 个体差异导致模型难以泛化
    \item 设备差异显著 —— 不同设备、导联、采样率带来分布偏移
    \item 模型泛化能力有限 —— 单任务训练，跨任务复用困难
\end{itemize}

\subsection{项目目标}

复现三个具有代表性的 EEG 解码模型，分析其设计思路、表征能力和泛化能力的演化：

\begin{center}
\Large\bfseries
EEGNet \quad $\longrightarrow$ \quad EEG Conformer \quad $\longrightarrow$ \quad CBraMod
\end{center}

% ================================================================
% 2. 模型复现
% ================================================================
\section{模型复现 (Methods)}
\subsection{模型结构比较}

\placeholderh{模型架构对比}{%
    EEGNet / EEG Conformer / CBraMod 结构示意图。\\
    展示三个模型的关键模块和架构差异。}{3.5cm}

\modelsizecompare

\textbf{主要发现：}从 EEGNet 到 EEG Conformer 再到 CBraMod，模型表现出更强的表征能力、更好的跨任务泛化能力，以及更复杂的模型结构。

\subsection{EEGNet · CNN 时代}

\textbf{核心思想：}利用 EEG 领域知识设计轻量化卷积网络（小参数量、强先验）。

\textbf{关键模块：}Temporal Convolution $\rightarrow$ Depthwise Convolution $\rightarrow$ Separable Convolution

\textbf{特点：}局部时空特征提取，参数量小，适合小样本 EEG 任务。

\subsection{EEG Conformer · CNN + Transformer 时代}

\textbf{核心思想：}结合 CNN 局部特征提取与 Transformer 全局建模。

\textbf{关键模块：}Convolution Module + Multi-Head Self-Attention

\textbf{特点：}同时建模局部特征与全局依赖，增强长时程信息捕获能力。

\subsection{CBraMod · Foundation Model 时代}

\textbf{核心思想：}基于大规模脑电数据预训练，学习通用 EEG 表征。

\textbf{关键模块：}EEG Tokenization $\rightarrow$ Criss-Cross Attention $\rightarrow$ Self-Supervised Pretraining

\textbf{特点：}学习通用脑电表征，支持跨任务迁移，提升泛化能力。

\vspace{0.5em}
{\small\centering
\textbf{三者关系：} EEGNet (CNN) $\rightarrow$ EEG Conformer (+Transformer) $\rightarrow$ CBraMod (+预训练)
\par}

% ================================================================
% 3. 复现结果
% ================================================================
\section{复现结果 (Results)}

\subsection{模型性能比较}

\placeholderh{模型性能对比}{%
    三个模型在多个数据集上的性能比较。\\
    指标：Accuracy / F1 Score / Cohen's Kappa。\quad (等待队友复现结果)}{4.5cm}

\subsection{训练效率比较}

\placeholderh{训练效率对比}{%
    三个模型的参数量 / 训练时间 / 推理速度对比。\quad (等待队友复现结果)}{3.2cm}

% ================================================================
% 4. 分析与讨论
% ================================================================
\section{分析与讨论 (Discussion)}

\subsection{观察一：模型能力持续提升}

\begin{center}
\large\bfseries
局部特征 \quad $\longrightarrow$ \quad 全局依赖 \quad $\longrightarrow$ \quad 通用表征
\end{center}

从 CNN 局部时空特征，到 Transformer 全局依赖，再到预训练通用表征——每一步都提升了表征空间和泛化边界。

\subsection{观察二：Foundation Model 成为新趋势}

CBraMod 通过大规模预训练学习通用 EEG 表示，支持多任务迁移，改变了 EEG 解码的训练范式：从单任务训练转向预训练后适配。

\subsection{观察三：可解释性成为新挑战}

Foundation Model 的 Embedding 是高维隐向量，虽然泛化更强，但生理意义难以解释——传统 Alpha/Beta/Theta 频段特征更容易与神经科学知识对应。

\begin{center}
\colorbox{myorange!10}{\parbox{0.9\columnwidth}{\centering\bfseries
    泛化能力 $\uparrow$ \hspace{2em} 但 \hspace{2em} 解释难度 $\uparrow$
}}
\end{center}

\placeholderh{可解释性挑战示意图}{%
    传统 EEG 分析 (可解释生理量) $\rightarrow$ Foundation Embeddings (高维隐空间)\\
    抽象程度递增，解释难度增大。\quad (等待绘制)}{2.5cm}

% ================================================================
% 5. 未来工作
% ================================================================
\section{未来工作 (Future Work)}

\subsection{NeuroGate：重新引入神经科学知识}

\begin{center}
\large\bfseries
Foundation Model \quad + \quad Neuroscience Prior
\end{center}

\textbf{思路：}CBraMod Embedding $+$ 神经科学先验（频带能量、时域统计等）$\rightarrow$ Gated Fusion $\rightarrow$ 分类预测。目标：在保持泛化能力的同时提高可解释性。

\placeholderh{NeuroGate 结构示意图}{%
    CBraMod Backbone + 神经科学先验 $\rightarrow$ Gated Fusion $\rightarrow$ 增强表征\\
    \small (等待绘制)}{2.8cm}

% ================================================================
% 6. 总结
% ================================================================
\section{总结 (Summary)}

\begin{center}
\colorbox{lightblue}{\parbox{0.95\columnwidth}{\centering
    {\Large\bfseries
        泛化能力持续提升\par}
    \vspace{0.8em}
    {\normalsize
        EEGNet/EEG Conformer $\rightarrow$ CBraMod\par}
    \vspace{0.5em}
    {\normalsize
        但可解释性问题逐渐凸显\par}
    \vspace{0.8em}
    {\large\bfseries
        未来方向：Foundation Model + Neuroscience Prior
        $\rightarrow$ 强泛化 + 强解释\par}
}}
\end{center}

% ================================================================
% 参考文献
% ================================================================
\vspace{0.5em}
\section*{参考文献}

{\footnotesize
\begin{enumerate}[leftmargin=1.5em, itemsep=0.1em]
    \item Lawhern et al. \textbf{EEGNet: A Compact CNN for EEG-based BCIs}. \textit{J. Neural Eng.}, 2018.
    \item Song et al. \textbf{EEG Conformer: Convolutional Transformer for EEG Decoding}. \textit{IEEE TNSRE}, 2023.
    \item Wang et al. \textbf{CBraMod: A Criss-Cross Brain Foundation Model for EEG Decoding}. \textit{ICLR}, 2025.
    \item NeuroGate / NeuroKE 内部参考：neural-prior gated fusion experiments.
\end{enumerate}
}

\end{multicols}

% =====================
% 底部
% =====================
\vspace{2mm}
{\color{myblue}\hrule height 1.5pt}
\vspace{2mm}
\begin{center}
{\footnotesize\color{mygray}
    BMEB215 · 机器学习与医学工程应用 · 2026 Spring \hspace{3em}
    南方科技大学 · 生物医学工程系
}
\end{center}

\end{document}
